% Part: incompleteness
% Chapter: incompleteness-provability
% Section: godels-paper

\documentclass[../../../include/open-logic-section]{subfiles}

\begin{document}

\olfileid{inc}{inp}{gop}

\olsection{गोडेलच्या मूळ शोधपत्राशी तुलना}

गोडेलच्या 1931 मधील शोधपत्राला थोडा वेळ
देणे उपयुक्त ठरेल. त्याच्या प्रस्तावनेत आपण
नुकत्याच पाहिलेल्या कल्पनांची रूपरेषा आहे.
शोधपत्र नुसते चाळले तरी प्रत्येक टप्प्यावर
काय चालले आहे हे सहज दिसते: प्रथम गोडेल
आकारिक पद्धत $P$ (विन्यास, स्वयंसिद्धके,
सिद्धतांचे नियम) वर्णन करतो; मग आदिम
पुनरावर्ती फलने आणि संबंध परिभाषित करतो;
त्यानंतर $x B y$ हा संबंध आदिम पुनरावर्ती
आहे हे दाखवतो आणि आदिम पुनरावर्ती फलने
व संबंध $\Th{P}$ मध्ये निरूपित होतात
असा युक्तिवाद करतो. मग वर दिल्याप्रमाणे
अपूर्णता प्रमेय सिद्ध करतो. तिसऱ्या विभागात,
सिद्ध करता न येणारे विधान अंकगणिताच्या
भाषेतील वाक्य म्हणून घेता येते हे दाखवतो.
यातूनच $\beta$-उपपत्ती उद्भवली; पुनरावर्ती
फलने $\Th{Q}$ मध्ये निरूपणीय आहेत हे
दाखवताना अनुक्रम हाताळण्यासाठी आपणही ती
वापरली. गोडेल एवढेच पुरेसे ठरणाऱ्या
किमान स्वयंसिद्धकसंचाची वेगळी ओळख करून
देत नाही, पण आता $\Th{Q}$ ते काम करते
हे आपल्याला माहीत आहे. अखेरीस चौथ्या
विभागात तो दुसऱ्या अपूर्णता प्रमेयाच्या
सिद्धतेची रूपरेषा देतो.

\end{document}
